\documentclass[10pt,a4paper]{article}
\usepackage[utf8]{inputenc}
\usepackage[vietnamese]{babel}
\usepackage{amsmath,amssymb}
\usepackage{geometry}
\geometry{
    a4paper,
    top=22mm,
    bottom=22mm,
    left=18mm,
    right=18mm,
    headheight=14pt,
    headsep=18pt,
    footskip=18pt
}
\usepackage{fancyhdr}
\usepackage{xcolor}
\usepackage{tcolorbox}
\tcbuselibrary{skins}
\usepackage{mathptmx}

\definecolor{stewartred}{RGB}{186, 12, 47}

\pagestyle{fancy}
\fancyhf{}
\renewcommand{\headrulewidth}{0pt}
\fancyhead[L]{\hspace*{0.29\textwidth}\fontsize{9pt}{11pt}\selectfont\textbf{\textsf{MỤC 3.1}}\quad\textsf{Đạo hàm của hàm đa thức và hàm mũ}}
\fancyhead[R]{\fontsize{9pt}{11pt}\selectfont\textbf{\textsf{175}}}

\newtcolorbox{definitionbox}[1][]{
    enhanced,
    colback=white,
    colframe=stewartred,
    arc=0mm,
    boxrule=1pt,
    left=10pt,
    right=10pt,
    top=7pt,
    bottom=7pt,
    sharp corners,
    #1
}

\setlength{\parindent}{0pt}
\setlength{\parskip}{6pt}

\begin{document}

\noindent
\begin{minipage}[t]{0.26\textwidth}
    \vspace*{318pt}
    \raggedright\fontsize{8.5pt}{11.5pt}\selectfont\textsf{Định lý Nhị thức được trình bày trên Trang Tra cứu 1.}
\end{minipage}%
\hfill
\begin{minipage}[t]{0.71\textwidth}
    \setlength{\parindent}{0pt}
    \setlength{\parskip}{6pt}
    
    So sánh các phương trình trong (1), (2) và (3), ta thấy xuất hiện một quy luật. Dự đoán rằng khi $n$ là một số nguyên dương thì $(d/dx)(x^n) = nx^{n-1}$ dường như là hợp lý. Điều này hóa ra hoàn toàn đúng. Chúng ta chứng minh quy tắc này theo hai cách; cách chứng minh thứ hai sử dụng Định lý Nhị thức.

    \vspace{2pt}
    \begin{definitionbox}
        \noindent\textbf{\textsf{\color{stewartred}Quy tắc lũy thừa}}\quad Nếu $n$ là một số nguyên dương, thì
        \vspace{-3pt}
        \[
        \frac{d}{dx}(x^n) = nx^{n-1}
        \]
    \end{definitionbox}
    \vspace{2pt}

    \noindent\textbf{\textsf{\color{stewartred}CÁCH CHỨNG MINH THỨ NHẤT}}\quad Công thức
    \[
    x^n - a^n = (x - a)(x^{n-1} + x^{n-2}a + \dots + xa^{n-2} + a^{n-1})
    \]
    có thể được kiểm chứng một cách đơn giản bằng cách nhân khai triển vế phải (hoặc bằng cách tính tổng thừa số thứ hai như một cấp số nhân). Nếu $f(x) = x^n$, chúng ta có thể sử dụng Phương trình 2.7.5 cho $f'(a)$ và phương trình trên để viết
    \[
    \begin{aligned}
    f'(a) = \lim_{x \to a} \frac{f(x) - f(a)}{x - a} &= \lim_{x \to a} \frac{x^n - a^n}{x - a} \\[5pt]
    &= \lim_{x \to a} (x^{n-1} + x^{n-2}a + \dots + xa^{n-2} + a^{n-1}) \\[5pt]
    &= a^{n-1} + a^{n-2}a + \dots + aa^{n-2} + a^{n-1} \\[5pt]
    &= na^{n-1}
    \end{aligned}
    \]

    \noindent\textbf{\textsf{\color{stewartred}CÁCH CHỨNG MINH THỨ HAI}}
    \[
    f'(x) = \lim_{h \to 0} \frac{f(x + h) - f(x)}{h} = \lim_{h \to 0} \frac{(x + h)^n - x^n}{h}
    \]

    Khi tìm đạo hàm của $x^4$, chúng ta đã phải khai triển $(x + h)^4$. Ở đây chúng ta cần khai triển $(x + h)^n$ và chúng ta sử dụng Định lý Nhị thức để làm điều đó:
    \begin{align*}
    f'(x) &= \lim_{h \to 0} \frac{\left[ x^n + nx^{n-1}h + \dfrac{n(n - 1)}{2}x^{n-2}h^2 + \dots + nxh^{n-1} + h^n \right] - x^n}{h} \\[7pt]
    &= \lim_{h \to 0} \frac{nx^{n-1}h + \dfrac{n(n - 1)}{2}x^{n-2}h^2 + \dots + nxh^{n-1} + h^n}{h} \\[7pt]
    &= \lim_{h \to 0} \left[ nx^{n-1} + \dfrac{n(n - 1)}{2}x^{n-2}h + \dots + nxh^{n-2} + h^{n-1} \right] \\[7pt]
    &= nx^{n-1}
    \end{align*}
    bởi vì mọi số hạng ngoại trừ số hạng đầu tiên đều có $h$ là một nhân tử và do đó tiến về 0.\hfill{\color{stewartred}\rule{1.2ex}{1.2ex}}

    \vspace{4pt}
    Chúng ta minh họa Quy tắc lũy thừa bằng nhiều ký hiệu khác nhau trong Ví dụ 1.
\end{minipage}

\end{document}